\section{Trabajos relacionados}

Se revisaron trabajos sobre reclutamiento universitario y selección de personal docente. Estos permiten comparar tanto el flujo de postulación como la evaluación de candidatos. Sus resultados sirven como referencia, pero no demuestran por sí solos que una arquitectura distribuida sea necesaria para el SSDC.

Enzelin et al. desarrollaron una aplicación de reclutamiento para una universidad mediante el método Rapid Application Development. La aplicación permite registrarse como postulante o empleador, cargar hojas de vida, consultar y publicar vacantes, y agendar entrevistas. Sus pruebas de aceptación reportaron satisfacción entre los usuarios \cite{enzelin2024recruitment}. Para el SSDC, este trabajo es pertinente por la carga de documentos y la gestión del proceso desde una aplicación. El proyecto de la UTEQ también debe representar la evaluación de méritos, el seguimiento de postulaciones y la preparación de resultados para las autoridades, según el alcance definido para esta entrega.

Dawami et al. aplicaron el método AHP como apoyo para seleccionar docentes en la Universitas Nahdlatul Ulama Sunan Giri. El estudio organiza criterios como formación reciente, promedio académico, logros, acreditación, experiencia y pruebas; luego calcula prioridades y compara el resultado con la selección realizada en la institución \cite{dawami2022ahp}. Su aporte para el SSDC es mostrar cómo una evaluación puede organizarse mediante criterios explícitos. Sin embargo, la matriz que utilizará la UTEQ debe basarse en sus propios criterios institucionales y no copiar automáticamente las ponderaciones de ese estudio.

Tadić et al. desarrollaron un modelo difuso TOPSIS para apoyar la selección de personal docente e investigador en educación superior. El modelo organiza criterios en una jerarquía y considera variables cualitativas, cuantitativas y las competencias requeridas de los expertos que evalúan \cite{tadic2022fuzzy}. Esta propuesta es relevante porque la selección académica puede incluir criterios de distinto tipo. El trabajo se centra en el modelo de decisión; el SSDC también debe gestionar las postulaciones, los documentos, los cambios de estado y la generación de resultados del proceso.

En conjunto, estos estudios aportan referencias para organizar la postulación digital y hacer explícitos los criterios de evaluación. El SSDC distribuido se diferencia porque busca representar el proceso de selección docente de la UTEQ y proponer servicios separados para sus funciones. Esa separación es una decisión arquitectónica del proyecto y deberá justificarse con los requisitos y las decisiones del equipo; no se atribuye como resultado de los estudios revisados.
